% 第4章小结。

\section{本章小结}
\label{sec:spp-summary}

章首的四类问题对应同一个过程的四个描述层次。固定一个时刻，过程是一个随机变量，回答单个时刻的读数可能是什么。固定有限个时刻，有限维分布记录它们的共同变化；相同边缘、不同相关的三个序列表明，平均值能否收敛、收敛多快都取决于这一层，而不取决于单时刻分布。固定一次实现，样本路径回答整条曲线的形状；在乘积$\sigma$代数中，全部有限维分布与路径规律是同一份信息，但其中每个事件只依赖可数个坐标，所以连续时间下“路径是否连续”“全程上确界多大”这类问题不能只靠有限维分布提出。最后，滤过规定每个时刻已经可以判定的事件：适应性要求当前值已知，停时要求停止决定不偷看未来。同一个过程相对于不同的滤过可以有不同的预测和不同的性质，所以鞅与Markov性质都必须连同滤过一起陈述；停时的定义也不保证停止后的期望等式，倍投策略说明需要另加条件。

构造过程有两条路线，分别对应有无时间顺序。有时间顺序时，初始分布加上依赖历史的概率核逐步给出有限路径的分布，核的归一化自动保证相容，Ionescu-Tulcea定理把它们拼成无穷序列的规律，且不需要拓扑条件；在标准Borel空间上，序列规律与“初始分布加条件核”是等价的描述。没有时间顺序时，给出全体有限子集上的相容分布族，Kolmogorov扩张定理在标准Borel状态空间上给出乘积空间上唯一的规律；两两相容而不存在三元规律的例子提醒我们，每个有限子集上都必须给出真正的概率测度。扩张定理交付的是坐标族，不是路径性质。版本与不可区分把“同一个过程”的两种含义分开：随机位置的尖点说明，不可数索引下二者确实不同，有限维分布不能决定路径是否连续；右连续性让可数个时刻代替全部时刻，连续版本因此在不可区分意义下唯一。Kolmogorov连续性准则用增量矩的幂次条件给出连续版本和Hölder指数，增量矩随间隔的标度把Brownian运动与Poisson计数过程区分开来。

后三节把索引推广到一般位置和集合。高斯过程的全部规律由均值函数与半正定协方差函数决定，半正定性既是必要条件也是充分条件；路径的正则性由增量方差的衰减控制，而同一个核作为再生核所定义的Hilbert空间描述的是确定性函数，Brownian路径几乎必然不在其中。交换序列的依赖来自共享的随机背景：交换性与独立、与平稳各约束联合规律的不同方面，de Finetti定理表明无限可交换序列就是“先抽随机分布、再条件独立抽样”，经验频率收敛到一个随机的极限。随机概率测度的每次实现本身是一个概率分布；有限分割上的相容规律不足以保证每次实现都可数可加，折棒构造给出了每次实现都是概率测度的Dirichlet过程，浓度参数与基测度分别控制离散程度与平均形状。

后两章按时间索引组织。第~\ref{chap:discrete-time-stochastic-processes}章在离散时间中使用本章的逐步核构造、滤过与停时，研究转移、累计、停止与长期行为；离散索引下版本与不可区分没有区别，乘积$\sigma$代数已经包含所需的路径事件。第~\ref{chap:continuous-time-stochastic-processes}章面对不可数时刻，需要本章的连续版本、联合可测性与路径正则性，并在相应问题处补充滤过的右连续性、可积性与过程列的统一控制等新条件。Markov性质、鞅性质与平稳性是可以同时成立的性质，不是互斥的过程类别。本章的高斯过程、交换序列与随机概率测度展示了构造的思路，但不是后两章每个结果的前提。
